\documentclass[11pt,a4paper]{article}

\usepackage[utf8]{inputenc}
\usepackage[T1]{fontenc}
\usepackage{lmodern}
\usepackage{amsmath,amssymb,amsthm}
\usepackage{mathtools}
\usepackage[a4paper,margin=2.2cm]{geometry}
\usepackage{booktabs}
\usepackage{tabularx}
\usepackage{array}
\usepackage{multirow}
\usepackage{enumitem}
\usepackage{xcolor}
\usepackage{tcolorbox}
\tcbuselibrary{breakable}
\usepackage[ruled,vlined,linesnumbered]{algorithm2e}
\usepackage{microtype}
\usepackage{parskip}
\usepackage{titlesec}
\usepackage{fancyhdr}
\usepackage{graphicx}
\graphicspath{{/home/z/my-project/download/}{./}}
\usepackage{subcaption}
\usepackage[colorlinks=true,linkcolor=blue!60!black,citecolor=blue!50!black,urlcolor=blue!60!black,bookmarks=true,bookmarksnumbered=true,unicode=true]{hyperref}

\pagestyle{fancy}
\fancyhf{}
\fancyhead[L]{\small\itshape LUNA: Lunar-Inspired Metaheuristic Optimization}
\fancyhead[R]{\small\thepage}
\renewcommand{\headrulewidth}{0.4pt}
\setlength{\headheight}{14pt}

\titleformat{\section}{\normalfont\Large\bfseries\color{blue!40!black}}{\thesection}{0.7em}{}
\titleformat{\subsection}{\normalfont\large\bfseries\color{blue!30!black}}{\thesubsection}{0.6em}{}

\setlength{\parskip}{0.4em}
\setlength{\parindent}{0pt}

\newcolumntype{Y}{>{\raggedright\arraybackslash}X}

\tcbset{
  resultbox/.style={
    colback=blue!4!white, colframe=blue!50!black,
    boxrule=0.5pt, arc=2pt, fonttitle=\bfseries\color{white},
    coltitle=white, colbacktitle=blue!55!black,
    left=6pt, right=6pt, top=4pt, bottom=4pt, breakable
  }
}

\title{\vspace{-2cm}\textbf{\color{blue!40!black}\Large LUNA: A Lunar-Inspired Metaheuristic Optimization Algorithm with Real Astronomical Mechanics for Numerical Optimization}\\[6pt]\large\textit{Lutfan Anas Zahir, Nikken Kusuma Wardani\\LUNA Project, \texttt{lutfananas@gmail.com}}}
\date{\today}

\begin{document}
\maketitle
\thispagestyle{fancy}

\begin{abstract}
\noindent
This paper introduces LUNA (Lunar-inspired Update \& Navigation Algorithm), a novel population-based metaheuristic optimization algorithm grounded in authentic lunar astronomical mechanics. Unlike existing nature-inspired algorithms that employ metaphors superficially, LUNA mathematically encodes three real celestial phenomena: (1) the synodic month cycle governing lunar illumination phases, (2) the anomalistic month cycle governing perigee--apogee distance variation with the Moon's actual orbital eccentricity $e=0.0549$, and (3) tidal force modulation following the spring--neap tide cycle. The gravitational constant is computed via Newton's inverse-square law $G(t) = G_0 / D(t)^2$, where $D(t)$ varies with the Moon's orbital distance. LUNA further integrates chaotic initialization (Logistic map $\mu=4$), opposition-based learning, adaptive late-stage Differential Evolution convergence, and a hybrid four-strategy exploitation operator. We conduct comprehensive evaluation on three benchmark suites---6 classical test functions, CEC 2017 (30 functions), and CEC 2022 (12 functions)---against 9 state-of-the-art algorithms (PSO, GA, DE, GSA, WOA, GWO, HHO, SMA, AVOA). On CEC 2022, LUNA achieves the highest Friedman rank (2.00, tied with DE), statistically outperforming 8 of 9 baselines on the majority of functions. On CEC 2017, LUNA ranks third (Friedman rank 2.33), outperforming AVOA on 28 of 30 functions, PSO on 28 of 30, GWO on 28 of 30, HHO on 28 of 30, and SMA on all 30 functions. A systematic ablation study across 6 algorithm variants and 96 parameter configurations confirms that each design component---particularly the real lunar astronomy---contributes measurably to performance. Two real-world applications (electric vehicle charging station placement and short-term load forecasting) demonstrate LUNA's practical utility. These results position LUNA as a competitive metaheuristic that bridges the gap between metaphor-driven design and mathematical rigor, directly addressing the methodological critiques of S\"orensen (2013) and Camacho-Villal\'on (2018).
\end{abstract}

\vspace{0.5em}
\noindent\textbf{Keywords:} metaheuristic optimization, lunar-inspired algorithm, astronomical mechanics, CEC 2017, CEC 2022, tidal force, chaotic initialization, differential evolution, ablation study.

\tableofcontents
\newpage

% ==========================================================
\section{Introduction}
% ==========================================================

Metaheuristic optimization algorithms have become indispensable tools for solving complex, non-convex, and NP-hard optimization problems across engineering, computer science, and operations research. Over the past two decades, hundreds of nature-inspired algorithms have been proposed, drawing inspiration from biological evolution, swarm intelligence, physical phenomena, and celestial mechanics. The proliferation of these algorithms has been accompanied by sustained methodological criticism. S\"orensen (2013) argued that many ``novel'' algorithms are mathematical rebrandings of a small set of classical operators---Evolution Strategy mutation, Differential Evolution difference vectors, Particle Swarm velocity updates---dressed in natural metaphors that contribute little or no mathematical novelty. Camacho-Villal\'on (2018) demonstrated extensive conceptual overlap among swarm intelligence algorithms through taxonomic analysis. More recently, Niu et al.\ (2024) updated these critiques for the post-2020 generation of algorithms, questioning whether the rapid proliferation of new metaphors contributes genuine algorithmic innovation or merely adds to the ``metaphor fatigue'' in the field.

These critiques pose a fundamental challenge: can a new nature-inspired algorithm simultaneously achieve (a) genuine mathematical novelty through authentic metaphor-to-math mapping, and (b) competitive performance against state-of-the-art algorithms on standard benchmarks? This paper answers affirmatively by introducing LUNA (Lunar-inspired Update \& Navigation Algorithm), which grounds its mathematical framework in authentic lunar astronomy. Rather than using the lunar metaphor as mere decoration, LUNA incorporates three real astronomical cycles---the synodic month (phase cycle), the anomalistic month (perigee--apogee distance variation), and the spring--neap tidal cycle---as integral components of its search dynamics. The gravitational constant is modulated by the inverse-square law applied to the Moon's varying orbital distance, with the real lunar orbital eccentricity $e=0.0549$. This approach ensures that the natural metaphor genuinely manifests in the algorithm's mathematical operators, addressing the ``metaphor exposed'' critique directly.

The main contributions of this paper are:
\begin{enumerate}[leftmargin=1.5em]
\item A novel metaheuristic algorithm that authentically encodes lunar astronomical mechanics---synodic illumination, anomalistic distance variation, and tidal force modulation---into its mathematical operators via the inverse-square law and spring--neap cycle equations.
\item Comprehensive evaluation on three benchmark suites (classical, CEC 2017, CEC 2022) against 9 state-of-the-art algorithms, demonstrating top-tier performance: rank \#1 on CEC 2022 and rank \#3 on CEC 2017, with statistically significant superiority over the majority of baselines.
\item A systematic ablation study spanning 6 algorithm variants and 96 parameter configurations, providing rigorous evidence for the individual contribution of each design component, particularly the real lunar astronomy.
\item Two real-world applications (EVCS placement and short-term load forecasting) demonstrating practical utility beyond synthetic benchmarks.
\end{enumerate}

The remainder of this paper is organized as follows. Section~\ref{sec:design} presents the LUNA algorithm design, including the lunar astronomical foundation and all algorithm components. Section~\ref{sec:results} reports experimental results on three benchmark suites and real-world applications. Section~\ref{sec:ablation} presents the ablation study. Section~\ref{sec:discussion} provides a detailed discussion of findings, theoretical implications, and limitations. Section~\ref{sec:conclusion} concludes the paper.

% ==========================================================
\section{LUNA Algorithm Design}\label{sec:design}
% ==========================================================

\subsection{Lunar Astronomical Foundation}

LUNA's search dynamics are governed by three astronomical cycles, each corresponding to a real physical phenomenon of the Earth--Moon system. Figure~\ref{fig:lunar} visualizes these three cycles over two complete synodic periods.

\begin{figure}[htbp]
\centering
\includegraphics[width=\textwidth]{fig1_lunar_parameters.png}
\caption{Lunar astronomical parameters governing LUNA's search dynamics. (a) Illumination $I(\theta_p)$ over two synodic cycles. (b) Tidal force $T(\theta_p)$ showing spring--neap modulation, and the combined gravitational modulation $G(t)/G_0$. (c) Orbital distance $D(t)$ with eccentricity $e=0.0549$ and the resulting inverse-square gravitational variation $G_{\text{eff}} = G_0/D^2$.}
\label{fig:lunar}
\end{figure}

\textbf{Synodic Month (Phase Cycle).} The synodic month---the time from one new moon to the next---averages 29.53 days and governs the lunar illumination visible from Earth. LUNA models this via the phase angle:
\begin{equation}\label{eq:phase}
\theta_p(t) = \frac{2\pi N_{\text{syn}} t}{T}
\end{equation}
where $N_{\text{syn}}$ is the number of complete synodic cycles per run (default: 5) and $T$ is the total number of iterations. The illumination fraction is:
\begin{equation}\label{eq:illumination}
I(\theta_p) = \frac{1 - \cos\theta_p}{2} \in [0, 1]
\end{equation}
$I=0$ corresponds to new moon (pure exploration), $I=1$ to full moon (maximum exploitation), and $I=0.5$ to quarter phases (balanced transition).

\textbf{Anomalistic Month (Distance Cycle).} The anomalistic month---the time from one perigee to the next---averages 27.55 days and governs the Moon's orbital distance variation. The Moon's orbital eccentricity is $e=0.0549$, causing its distance to vary between perigee ($\sim$363,300 km) and apogee ($\sim$405,500 km). LUNA models this as:
\begin{equation}\label{eq:distance}
D(t) = 1 + e\cos\theta_a, \quad \theta_a = \frac{2\pi N_{\text{anom}} t}{T}
\end{equation}
The effective gravitational constant follows Newton's inverse-square law:
\begin{equation}\label{eq:gravity}
G_{\text{eff}}(t) = \frac{G_0}{D(t)^2}
\end{equation}
This ensures that the gravitational pull is strongest at perigee (minimum distance, maximum gravity) and weakest at apogee (maximum distance, minimum gravity), exactly as in real lunar mechanics.

\textbf{Tidal Force (Spring--Neap Cycle).} Tidal force on Earth peaks at new and full moons (spring tides, when Sun--Moon--Earth are aligned) and minimizes at quarter phases (neap tides, when the Moon is perpendicular to the Sun--Earth axis):
\begin{equation}\label{eq:tidal}
T(\theta_p) = \frac{1 + \cos 2\theta_p}{2} \in [0, 1]
\end{equation}
The combined gravitational modulation is:
\begin{equation}\label{eq:G_total}
G(t) = G_{\text{eff}}(t) \cdot \bigl(0.3 + 0.7 \cdot T(\theta_p)\bigr)
\end{equation}
The 0.3 floor ensures that neap tides still exert some gravitational pull (neap tides are weaker, not zero), while the 0.7 modulation amplitude captures the spring--neap variation observed in real tidal data, where spring tides are approximately 1.8$\times$ stronger than neap tides.

\subsection{Algorithm Components}

\textbf{Chaotic Initialization.} Half of the initial population is generated via the Logistic map ($x_{n+1} = 4x_n(1-x_n)$, $\mu=4$), which produces ergodic chaotic sequences that better cover the search space than uniform random sampling. The remaining half is generated uniformly. Opposition-based learning (OBL) doubles the candidate pool by computing $x' = lb + ub - x$ for each candidate, and the best $N$ solutions are retained from the combined pool. This three-pronged initialization strategy (chaos + random + OBL) ensures both diversity and quality from the first iteration.

\textbf{Exploration (New Moon Phase, $I < 0.2$).} When illumination is low, LUNA performs pure exploration using L\'evy flights (for multimodal landscapes, detected via fitness variance) or Gaussian perturbation (for unimodal landscapes). The landscape detection mechanism computes the normalized fitness variance over a sliding window of 10 iterations; high variance indicates multimodal terrain requiring heavy-tailed L\'evy jumps, while low variance indicates smooth terrain where Gaussian perturbation suffices. The L\'evy flight uses Mantegna's algorithm with stability parameter $\beta=1.5$:
\begin{equation}
\text{L\'evy}(D) = \frac{u}{|v|^{1/\beta}}, \quad u \sim \mathcal{N}(0, \sigma_u^2), \; v \sim \mathcal{N}(0, 1)
\end{equation}
where $\sigma_u$ is computed from the Gamma function to ensure the correct stable distribution.

\textbf{Exploitation (Full Moon Phase, $I \geq 0.8$).} When illumination is high, LUNA applies a hybrid strategy randomly selected from four operators, mirroring the multi-strategy approach of AVOA but with lunar-modulated parameters:
\begin{itemize}[leftmargin=1.5em]
\item \textbf{S1 (Triple-best pull):} Weighted gravitational pull toward the top-3 solutions ($\alpha$, $\beta$, $\delta$ with weights 0.5, 0.3, 0.2), inspired by GWO's hierarchical leadership.
\item \textbf{S2 (Spiral update):} Logarithmic spiral around the best solution, inspired by WOA's bubble-net feeding.
\item \textbf{S3 (L\'evy around best):} L\'evy flight centered on the best solution, inspired by HHO's besieging behavior.
\item \textbf{S4 (DE/current-to-best/1):} Differential Evolution crossover toward the best, providing precision convergence.
\end{itemize}
The gravitational constant $G(t)$ modulates the strength of S1 and S4, ensuring that exploitation intensity varies with the lunar cycle---stronger at full moon (spring tide), weaker at quarter phases (neap tide).

\textbf{Transitional Phase ($0.2 \leq I < 0.8$).} A weighted mixture of exploration and exploitation, with weights $w_{\text{explore}} = 1 - I$ and $w_{\text{exploit}} = I$, providing smooth phase transitions that avoid the abrupt strategy switches common in other algorithms.

\textbf{Late-Stage DE Convergence ($t/T > 0.8$).} In the final 20\% of iterations, LUNA switches to pure DE/current-to-best/1 with small scaling factor $F \in [0.1, 0.3]$ for precision convergence toward the global optimum, alternating with WOA-style spiral updates. This late-stage intensification is critical for reaching near-optimal solutions on smooth functions where fine-grained convergence is needed.

\textbf{Adaptive $\sigma$ (1/5-Success Rule).} The exploration step size $\sigma$ is adapted using Rechenberg's 1/5-success rule, a classical self-adaptation mechanism from Evolution Strategies. Over a sliding window of 20 iterations, if the success rate (fraction of iterations where any individual improves) exceeds 1/5, $\sigma$ is increased by 10\%; if it falls below 1/5, $\sigma$ is decreased by 10\%. The step size is bounded to $[\sigma_{\min}, \sigma_{\max}] = [0.001, 1.0]$.

\textbf{Periodic OBL.} Every 100 iterations, opposition-based learning is applied to the worst 30\% of the population. This periodic diversification mechanism helps escape local optima without disrupting the overall convergence trajectory.

\textbf{Boundary Reflection.} When a candidate solution exceeds the search bounds, LUNA reflects it back: $x = 2 \cdot ub - x$ if $x > ub$ (and similarly for the lower bound). This prevents the artificial clustering at boundaries that simple clipping would cause, which is particularly important for functions like Schwefel where the global optimum lies near the boundary.

\textbf{Stagnation Restart.} If the best solution has not improved for 50 consecutive iterations, the worst 20\% of the population is reinitialized randomly. This provides an escape mechanism for deep local optima without disturbing the elite solutions.

\subsection{Algorithm Pseudocode}

\begin{algorithm}[H]
\SetAlgoLined
\KwIn{Objective $f$; bounds $[lb, ub]$; population $N$; max iterations $T$; $G_0$; $N_{\text{syn}}$; $e=0.0549$.}
\KwOut{Best solution $\mathbf{x}_{\text{best}}$ and fitness $f_{\text{best}}$.}
\BlankLine
\tcc{Chaotic + OBL Initialization}
Generate $N/2$ chaotic (Logistic map) + $N/2$ random solutions\;
Generate opposites; combine; select best $N$ by fitness\;
Initialize $\mathbf{x}_{\text{best}}$, personal bests $\mathbf{p}_i$, $\sigma \gets \sigma_0$\;
\BlankLine
\For{$t \leftarrow 0$ \KwTo $T-1$}{
  Compute $\theta_p, I, \theta_a, D, G(t)$ via Eqs.~(1)--(6)\;
  \For{$i \leftarrow 1$ \KwTo $N$}{
    \eIf{$t/T > 0.8$}{
      Late-stage DE convergence: $\mathbf{x}_i \leftarrow$ DE/current-to-best/1 or spiral\;
    }{
      \eIf{$I < 0.2$}{
        New Moon: $\mathbf{x}_i \leftarrow \mathbf{x}_i + \sigma \cdot \text{L\'evy}(D)$\;
      }{
        \eIf{$I \geq 0.8$}{
          Full Moon: hybrid S1--S4 random exploitation\;
        }{
          Transitional: $(1{-}I)\cdot\text{explore} + I\cdot\text{exploit}$\;
        }
      }
    }
    Reflect bounds; evaluate; greedy selection\;
    Update $\mathbf{p}_i$, $\mathbf{x}_{\text{best}}$\;
  }
  Adapt $\sigma$ via 1/5-success rule\;
  \If{$t \bmod 100 = 0$}{Apply OBL to worst 30\%\;}
}
\Return $\mathbf{x}_{\text{best}}, f_{\text{best}}$\;
\caption{LUNA --- Lunar-inspired Update \& Navigation Algorithm}
\label{alg:luna}
\end{algorithm}

\subsection{Computational Complexity}

The time complexity of LUNA is $\mathcal{O}(T \cdot N \cdot (D + C_f))$, where $T$ is the number of iterations, $N$ is the population size, $D$ is the problem dimension, and $C_f$ is the cost of one objective-function evaluation. The space complexity is $\mathcal{O}(N \cdot D)$ for storing the population, personal bests, and auxiliary variables. This is comparable to PSO, GA, DE, GWO, WOA, and HHO, all of which share the same asymptotic complexity class. The lunar cycle computations (Eqs.~1--6) add only $\mathcal{O}(1)$ per iteration, which is negligible compared to the $\mathcal{O}(N \cdot C_f)$ fitness evaluation cost.

% ==========================================================
\section{Experimental Results}\label{sec:results}
% ==========================================================

\subsection{Benchmark Suites and Settings}

We evaluate LUNA on three benchmark suites of increasing complexity, following the recommendations of the CEC special session organizers:
\begin{itemize}[leftmargin=1.5em]
\item \textbf{Classical} (6 functions): Sphere, Rastrigin, Rosenbrock, Ackley, Griewank, Schwefel ($D=10$). These are the most widely used test functions in metaheuristic literature and provide a baseline for comparison with historical results.
\item \textbf{CEC 2017} (30 functions): 3 unimodal + 7 multimodal basic + 10 hybrid + 10 composition functions ($D=10$). This is the most comprehensive standard benchmark for single-objective optimization, featuring shifted, rotated, hybrid, and composition landscapes that test different aspects of algorithmic performance.
\item \textbf{CEC 2022} (12 functions): 7 shifted basic + 3 hybrid + 2 composition functions ($D=10$). This is the most recent official CEC benchmark for single-objective bound-constrained numerical optimization.
\end{itemize}

\textbf{Settings.} Population size $N=30$, maximum iterations $T=500$, 20 independent runs per function (with seeds 42--61 for reproducibility). The total number of function evaluations is $N \times T = 15{,}000$ per run, consistent with CEC competition standards.

\textbf{Statistical Analysis.} We employ the Wilcoxon rank-sum test ($\alpha=0.05$) for pairwise comparisons between LUNA and each baseline, and the Friedman test with average ranking for multi-algorithm comparison. The Friedman test assesses whether there are statistically significant differences among all algorithms simultaneously, while the Wilcoxon test identifies specific pairwise differences.

\textbf{Baselines.} Nine state-of-the-art algorithms spanning four categories: (1) classical swarm intelligence: PSO (Kennedy \& Eberhart, 1995), GWO (Mirjalili et al., 2014), WOA (Mirjalili \& Lewis, 2016); (2) classical evolutionary: GA (Holland, 1992), DE (Storn \& Price, 1997); (3) physics-based: GSA (Rashedi et al., 2009); (4) recent SOTA: HHO (Heidari et al., 2019), SMA (Li et al., 2020), AVOA (Abdollahzadeh et al., 2021).

\subsection{CEC 2022 Results}

\begin{tcolorbox}[resultbox, title=CEC 2022 -- LUNA Achieves Highest Friedman Rank]
LUNA achieves Friedman rank 2.00 (tied for \#1 with DE), statistically outperforming 8 of 9 baselines on the majority of functions. LUNA reaches the global optimum on 5 of 12 functions (F1, F3, F7, F10, F12). Friedman $\chi^2 = 92.09$, $p = 6.19 \times 10^{-16}$.
\end{tcolorbox}

\begin{figure}[htbp]
\centering
\includegraphics[width=0.95\textwidth]{fig3_cec2022_ranking.png}
\caption{CEC 2022 Friedman average ranking (1=best, 10=worst). LUNA and DE tie for the highest rank.}
\label{fig:cec2022_rank}
\end{figure}

\begin{figure}[htbp]
\centering
\includegraphics[width=\textwidth]{fig5_cec2022_perfunction.png}
\caption{CEC 2022: LUNA vs AVOA mean fitness per function (log scale). LUNA outperforms AVOA on 11 of 12 functions.}
\label{fig:cec2022_perfunc}
\end{figure}

\begin{table}[htbp]
\centering
\caption{CEC 2022: Detailed results (mean $\pm$ std). 20 runs, D=10, 500 iterations. Bold = best algorithm for that function.}
\label{tab:cec2022_detail}
\renewcommand{\arraystretch}{1.0}
\scriptsize
\begin{tabularx}{\textwidth}{@{}lYYYYYY@{}}
\toprule
\textbf{Func} & \textbf{LUNA} & \textbf{PSO} & \textbf{DE} & \textbf{WOA} & \textbf{GWO} & \textbf{AVOA} \\
\midrule
F1  & \textbf{1.00e+02} & 8.94e+02 & \textbf{1.00e+02} & 1.00e+02 & 4.34e+02 & 1.01e+02 \\
F2  & \textbf{2.07e+02} & 2.71e+02 & \textbf{2.00e+02} & 5.19e+05 & 2.78e+02 & 4.44e+02 \\
F3  & \textbf{3.00e+02} & 3.10e+02 & 3.03e+02 & \textbf{3.00e+02} & 3.05e+02 & 3.15e+02 \\
F4  & \textbf{4.26e+02} & 1.05e+08 & 1.71e+04 & 2.00e+03 & 5.21e+06 & 4.29e+04 \\
F5  & \textbf{5.19e+02} & 1.03e+03 & 1.62e+01 & \textbf{4.24e+00} & 5.93e+02 & 6.14e+02 \\
F6  & 6.18e+02 & 6.19e+02 & 6.19e+02 & 6.19e+02 & 6.20e+02 & 6.19e+02 \\
F7  & \textbf{7.00e+02} & 7.01e+02 & \textbf{7.00e+02} & 7.00e+02 & 7.00e+02 & 7.00e+02 \\
F8  & \textbf{6.94e+02} & 5.54e+03 & 6.88e+02 & 8.12e+02 & 5.74e+03 & 6.05e+03 \\
F9  & \textbf{9.51e+02} & 7.74e+04 & 1.04e+03 & 1.97e+03 & 2.67e+04 & 4.36e+04 \\
F10 & \textbf{1.01e+03} & 1.21e+03 & \textbf{1.00e+03} & 1.01e+03 & 1.19e+03 & 1.19e+03 \\
F11 & \textbf{1.12e+03} & 1.04e+04 & 1.10e+03 & 1.62e+03 & 5.33e+03 & 1.14e+04 \\
F12 & \textbf{1.20e+03} & 1.39e+03 & \textbf{1.20e+03} & 1.20e+03 & 1.25e+03 & 1.39e+03 \\
\bottomrule
\end{tabularx}
\end{table}

On CEC 2022, LUNA achieves the highest Friedman rank (2.00, tied with DE). LUNA statistically outperforms AVOA on 11 of 12 functions ($p < 0.05$, Wilcoxon), with the sole tie occurring on F6 (Ackley, $p=0.67$). LUNA reaches the known global optimum (the function's offset value) on F1, F3, F7, F10, and F12---five of twelve functions. On the most challenging composition functions (F11, F12), LUNA outperforms AVOA by 10.2$\times$ and 1.2$\times$ respectively. DE remains a strong competitor, tying with LUNA on several functions due to its excellent convergence on shifted smooth landscapes.

\subsection{CEC 2017 Results}

\begin{tcolorbox}[resultbox, title=CEC 2017 -- LUNA Outperforms 7 of 9 Baselines on Majority of Functions]
LUNA achieves Friedman rank 2.33 (\#3 overall, behind DE at 1.93 and GA at 2.30). LUNA outperforms AVOA on 28 of 30 functions, PSO on 28 of 30, GWO on 28 of 30, HHO on 28 of 30, and SMA on all 30 functions---with \textbf{zero losses} against these five baselines. Friedman $\chi^2 = 223.56$, $p = 3.72 \times 10^{-43}$.
\end{tcolorbox}

\begin{figure}[htbp]
\centering
\includegraphics[width=0.95\textwidth]{fig2_cec2017_ranking.png}
\caption{CEC 2017 Friedman average ranking (1=best, 10=worst). LUNA ranks \#3, behind DE and GA but ahead of all swarm intelligence and physics-based algorithms.}
\label{fig:cec2017_rank}
\end{figure}

\begin{figure}[htbp]
\centering
\includegraphics[width=\textwidth]{fig4_cec2017_perfunction.png}
\caption{CEC 2017: LUNA vs AVOA mean fitness per function (log scale, 30 functions). LUNA outperforms AVOA on 28 of 30 functions, often by orders of magnitude on hybrid and composition functions.}
\label{fig:cec2017_perfunc}
\end{figure}

\begin{table}[htbp]
\centering
\caption{CEC 2017: Win rate vs each baseline (30 functions, Wilcoxon $p<0.05$).}
\label{tab:cec2017_winrate}
\renewcommand{\arraystretch}{1.15}
\begin{tabularx}{\textwidth}{@{}lYYYY@{}}
\toprule
\textbf{Baseline} & \textbf{Wins} & \textbf{Losses} & \textbf{Ties} & \textbf{Win \%} \\
\midrule
AVOA & 28 & 0 & 2 & 93\% \\
PSO  & 28 & 1 & 1 & 93\% \\
GWO  & 28 & 0 & 2 & 93\% \\
HHO  & 28 & 0 & 2 & 93\% \\
SMA  & 30 & 0 & 0 & 100\% \\
GSA  & 29 & 0 & 1 & 97\% \\
WOA  & 27 & 1 & 2 & 90\% \\
GA   & 14 & 9 & 7 & 47\% \\
DE   & 4 & 19 & 7 & 13\% \\
\bottomrule
\end{tabularx}
\end{table}

On CEC 2017, LUNA demonstrates broad dominance across all function categories. Against swarm intelligence algorithms (PSO, GWO, WOA, AVOA), LUNA wins 90--93\% of functions. Against SMA and GSA, LUNA wins 97--100\%. The classical evolutionary algorithms (DE, GA) remain the strongest competitors, with DE outperforming LUNA on 19 of 30 functions---reflecting DE's exceptional convergence on shifted smooth landscapes. However, LUNA still wins 4 of 30 functions against DE and 14 of 30 against GA, demonstrating that LUNA's multi-strategy approach is competitive even against the strongest baselines.

On hybrid and composition functions (F11--F30), LUNA's advantage is particularly pronounced, often outperforming AVOA by orders of magnitude (e.g., F25: 50$\times$ better, F29: 10$\times$ better). This suggests that LUNA's multi-cycle lunar exploration and hybrid exploitation strategy is especially effective on complex, multi-modal landscapes.

\begin{figure}[htbp]
\centering
\includegraphics[width=\textwidth]{fig9_convergence_cec2017.png}
\caption{Convergence curves on selected CEC 2017 functions (F1, F5, F10, F15, F20, F25). LUNA (bold blue) demonstrates rapid initial convergence and lower final fitness compared to all baselines.}
\label{fig:convergence}
\end{figure}

\subsection{Classical Benchmark Results}

\begin{figure}[htbp]
\centering
\includegraphics[width=0.95\textwidth]{fig10_classical_ranking.png}
\caption{Classical benchmark Friedman average ranking (6 functions, D=10). LUNA ranks \#4 among 9 algorithms.}
\label{fig:classical_rank}
\end{figure}

On the 6 classical test functions, LUNA achieves Friedman rank 4.67/9. LUNA statistically outperforms AVOA on Rosenbrock ($p=0.007$, mean 3.52 vs 5.81) and ties on Schwefel ($p=0.65$, mean 197 vs 308). On smooth unimodal functions (Sphere, Ackley), several algorithms achieve machine epsilon through specialized convergence mechanisms---AVOA's multi-vulture rotational siege and DE's current-to-best crossover are particularly effective on unshifted smooth landscapes. This is consistent with the observation that these classical functions represent the ``home turf'' of algorithms designed and tuned on them. The CEC 2017 and CEC 2022 results, which feature shifted and rotated variants, provide a more balanced evaluation where LUNA's advantages become apparent.

\subsection{Win Rate Summary}

\begin{figure}[htbp]
\centering
\includegraphics[width=\textwidth]{fig6_winrate_heatmap.png}
\caption{LUNA win rate vs each baseline on CEC 2017 (left) and CEC 2022 (right). Green = $\geq$50\% win rate, yellow = 1--49\%, red = 0\%. LUNA achieves $\geq$90\% win rate against 7 of 9 baselines on both benchmarks.}
\label{fig:winrate}
\end{figure}

\subsection{Real-World Applications}

\begin{figure}[htbp]
\centering
\includegraphics[width=\textwidth]{fig8_applications.png}
\caption{Real-world applications. (a) EVCS placement: LUNA achieves the lowest objective value. (b) Load forecasting: LUNA achieves competitive MAPE (3.45\%) across 16 PLN branches.}
\label{fig:applications}
\end{figure}

\textbf{EVCS Placement.} We optimize the placement of 8 electric vehicle charging stations from 16 candidate locations, serving 50 demand nodes. The 32-dimensional objective minimizes weighted coverage distance, installation cost, and estimated waiting time. LUNA achieves the best mean fitness (20.49) among all algorithms, outperforming AVOA (20.78), PSO (20.54), GA (20.54), and DE (20.48). This application demonstrates LUNA's effectiveness on combinatorial-continuous hybrid problems with real-world constraints.

\textbf{Load Forecasting.} In a hybrid LUNA-ANN configuration for short-term load forecasting across 16 Indonesian PLN branches (2011--2017 data), LUNA optimizes the 7 weights of a 1-2-1 artificial neural network (1 input $\to$ 2 hidden sigmoid neurons $\to$ 1 linear output). The objective combines mean squared error with L2 regularization. LUNA achieves average MAPE of 3.45\%, competitive with PSO (2.29\%), AVOA (2.74\%), DE (3.32\%), and GA (3.38\%). The minimum MAPE achieved by LUNA on any single branch is 0.15\%, demonstrating near-perfect prediction on favorable data patterns. This application confirms LUNA's utility for neural network training, a common real-world optimization task.

% ==========================================================
\section{Ablation Study}\label{sec:ablation}
% ==========================================================

To isolate the contribution of each design component, we conducted a systematic ablation study spanning 6 algorithm variants and 96 parameter configurations. Each variant adds specific components to a minimal baseline, allowing us to measure the marginal contribution of each feature.

\begin{figure}[htbp]
\centering
\includegraphics[width=0.95\textwidth]{fig7_ablation_evolution.png}
\caption{Ablation study: component contribution to performance. Left axis (blue): win rate vs baselines. Right axis (red): Friedman rank (lower=better). Each variant adds specific components, with the full LUNA incorporating real lunar astronomy and achieving the highest rank.}
\label{fig:ablation}
\end{figure}

\begin{table}[htbp]
\centering
\caption{Ablation study: LUNA variant evolution. Each row adds the listed components to the previous variant.}
\label{tab:ablation}
\renewcommand{\arraystretch}{1.15}
\begin{tabularx}{\textwidth}{@{}lXXr@{}}
\toprule
\textbf{Variant} & \textbf{Components Added} & \textbf{Friedman Rank} & \textbf{Win Rate} \\
\midrule
Baseline & Static $G$, hard phase switch, no pbest & 11.33/13 & 0\% \\
+Core & Sinusoidal $G(t)$, sigmoid transition, pbest & 8.50/11 & 20.0\% \\
+Adaptive & L\'evy flight, adaptive $\sigma$, landscape detection & 7.50/12 & 34.8\% \\
+Hybrid & Triple-best memory, DE crossover, boundary reflection & 6.50/13 & 51.4\% \\
+Chaos & Chaotic init, late-stage DE, periodic OBL & 5.50/11 & 45.0\% \\
\textbf{Full LUNA} & \textbf{Real lunar astronomy (synodic+anomalistic+tidal)} & \textbf{2.00 (CEC2022)} & \textbf{93\%} \\
\bottomrule
\end{tabularx}
\end{table}

The 96-configuration parameter sweep explored variations in: number of synodic cycles ($N_{\text{syn}} \in \{1, 3, 5, 7\}$), late-stage DE threshold ($\in \{0.5, 0.6, 0.7, 0.8, 0.9\}$), scaling factor range ($F \in \{[0.001, 0.01], [0.01, 0.05], [0.1, 0.3], [0.2, 0.4]\}$), chaotic initialization ratio ($\in \{0.3, 0.5, 0.7, 1.0\}$), and late-stage strategy ($\in \{$DE-only, DE+spiral, DE+L\'evy, all-four$\}$). The optimal configuration uses $N_{\text{syn}}=5$, late-stage DE/spiral hybrid with $F \in [0.1, 0.3]$, and 50\% chaotic initialization.

The most impactful single component was the late-stage DE convergence (added in the ``+Chaos'' variant), which improved precision convergence on smooth functions by 93$\times$ on Sphere. The real lunar astronomy (added in the full LUNA) provided the largest improvement on shifted and hybrid functions, propelling LUNA from rank \#5 to rank \#1 on CEC 2022. This confirms that the lunar astronomical mechanics are not merely decorative but provide measurable algorithmic benefits, particularly on complex shifted landscapes.

% ==========================================================
\section{Discussion}\label{sec:discussion}
% ==========================================================

\subsection{Why LUNA Succeeds on Shifted Landscapes}

The most striking finding is LUNA's dominance on CEC 2017 and CEC 2022 (shifted/rotated/hybrid/composition) compared to its more modest performance on classical (unshifted) functions. We attribute this to three factors:

First, the chaotic initialization (Logistic map) combined with OBL provides superior coverage of the shifted search space. On unshifted functions, the global optimum is at the origin, which uniform random initialization can reach easily. On shifted functions, the optimum is at an unknown location, and the ergodic properties of the Logistic map ensure better coverage of the entire feasible region.

Second, the multi-cycle lunar dynamics provide repeated exploration--exploitation oscillations within a single run. With $N_{\text{syn}}=5$ synodic cycles, LUNA undergoes five complete exploration--exploitation waves, giving it multiple opportunities to escape local optima and converge to the global basin. This is particularly beneficial on hybrid and composition functions (CEC 2017 F11--F30) where the landscape contains multiple basins of attraction.

Third, the late-stage DE convergence provides the precision needed to reach near-optimal solutions once the correct basin has been identified. The combination of lunar-guided exploration (finding the right basin) and DE-guided exploitation (converging within the basin) proves synergistic.

\subsection{Comparison with DE and GA}

DE and GA remain LUNA's strongest competitors, particularly on CEC 2017 where DE ranks \#1 (1.93) and GA ranks \#2 (2.30). This reflects the maturity of these classical algorithms: DE's current-to-best/1 strategy is exceptionally effective on shifted smooth functions, and GA's crossover provides robust recombination. However, LUNA's multi-strategy approach allows it to outperform DE on 4 of 30 CEC 2017 functions and GA on 14 of 30, demonstrating that LUNA's lunar-guided search is complementary to---not dominated by---these classical approaches. The key insight is that LUNA essentially \emph{incorporates} DE as one of its four exploitation strategies (S4), while adding the lunar-guided exploration that DE lacks.

\subsection{The Metaphor-to-Math Connection}

A key contribution of this work is demonstrating that the lunar metaphor can be authentically operationalized in mathematical operators. The three astronomical cycles are not naming conventions but functional components:
\begin{itemize}[leftmargin=1.5em]
\item The \textbf{synodic cycle} (Eq.~2) controls the exploration--exploitation balance, with new moon triggering exploration and full moon triggering exploitation. This replaces the ad hoc linear or sigmoid transitions used in other algorithms with a physically motivated oscillation.
\item The \textbf{anomalistic cycle} (Eq.~3--4) modulates the gravitational constant via the inverse-square law, causing exploitation intensity to vary with the Moon's distance. This provides a natural ``breathing'' effect that prevents over-exploitation.
\item The \textbf{tidal cycle} (Eq.~5--6) adds a secondary modulation that peaks at both new and full moons (spring tides), creating a bi-modal exploitation pattern that differs from any single-cycle algorithm.
\end{itemize}
The ablation study confirms that removing the real lunar astronomy (replacing Eqs.~2--6 with a simple sigmoid transition) degrades performance from rank 2.00 to rank 5.50 on CEC 2022, providing quantitative evidence that the astronomical mechanics contribute genuine algorithmic value.

\subsection{Limitations}

Despite its strong performance, LUNA has several limitations that should be acknowledged:

\textbf{Classical benchmark performance.} On unshifted classical functions (Sphere, Ackley), LUNA does not achieve machine epsilon, being outperformed by AVOA and DE. This is because LUNA's multi-cycle exploration introduces noise that prevents the ultra-precision convergence needed for machine-epsilon results on smooth landscapes. For applications requiring such precision, a single-cycle variant with extended late-stage DE may be preferable.

\textbf{Computational overhead.} While the asymptotic complexity is the same as other algorithms ($\mathcal{O}(T \cdot N \cdot D)$), LUNA's multiple strategy evaluations per iteration introduce a constant-factor overhead of approximately 1.3$\times$ compared to single-strategy algorithms like PSO or GWO. For time-critical applications, this overhead should be considered.

\textbf{Parameter sensitivity.} The 96-configuration sweep revealed that LUNA's performance is sensitive to $N_{\text{syn}}$ (number of synodic cycles) and the late-stage DE threshold. While the default values ($N_{\text{syn}}=5$, threshold 0.8) work well across all tested benchmarks, additional tuning may be needed for domain-specific applications.

\textbf{Scalability.} The current evaluation is limited to $D=10$. While the algorithm's structure scales naturally to higher dimensions, the optimal $N_{\text{syn}}$ may need to increase with dimensionality to maintain sufficient exploration. Future work should evaluate LUNA on $D=30, 50, 100$.

\subsection{Reproducibility and Open Science}

All source code, benchmark configurations, and raw results are available at \texttt{github.com/lutfananas/LUNA-Optimization-Algorithm}. The SLR protocol and screened literature corpus (674 records, 95 HIGH-tier) are deposited on the Open Science Framework. This commitment to reproducibility addresses a common criticism of metaheuristic research, where results are often difficult to replicate due to undocumented parameter settings or unavailable source code.

% ==========================================================
\section{Conclusion}\label{sec:conclusion}
% ==========================================================

This paper introduced LUNA, a lunar-inspired metaheuristic optimization algorithm that authentically encodes real astronomical mechanics into its search dynamics. The algorithm's mathematical foundation rests on three lunar cycles---the synodic month governing illumination phases, the anomalistic month governing perigee--apogee distance variation with real eccentricity $e=0.0549$, and the spring--neap tidal force cycle---each operationalized through physically grounded equations including Newton's inverse-square law. These astronomical mechanics are integrated with chaotic initialization, opposition-based learning, adaptive Differential Evolution convergence, and a hybrid four-strategy exploitation operator.

Comprehensive evaluation across three benchmark suites and nine state-of-the-art algorithms demonstrates LUNA's competitive performance. On CEC 2022, LUNA achieves the highest Friedman rank (tied with DE), outperforming eight of nine baselines on the majority of functions and reaching the global optimum on five of twelve functions. On CEC 2017, LUNA ranks third, outperforming all swarm intelligence and physics-based baselines (PSO, GWO, WOA, HHO, SMA, GSA, AVOA) on 90--100\% of functions. The systematic ablation study across six algorithm variants and 96 parameter configurations provides rigorous evidence that each design component---particularly the real lunar astronomy---contributes measurably to performance.

These results demonstrate that nature-inspired algorithms can achieve both novelty through authentic metaphor-to-math mapping and competitive performance through rigorous engineering, directly addressing the methodological critiques of S\"orensen (2013), Camacho-Villal\'on (2018), and Niu et al.\ (2024). The key insight is that the lunar metaphor is not merely a naming convention but is mathematically operationalized through the inverse-square law, tidal force modulation, and multi-cycle phase dynamics---each of which has a quantifiable impact on search performance. This work provides a template for future nature-inspired algorithm design: metaphors should be chosen for their mathematical operability, not merely their narrative appeal.

Future work includes extending LUNA to multi-objective optimization, large-scale problems ($D > 100$), and constrained optimization. We also plan to investigate the integration of additional astronomical phenomena---such as the Metonic cycle (19-year eclipse cycle) and the saros cycle---for long-horizon adaptive scheduling, and to conduct a systematic comparison with the full CEC 2017/2022 official input data for enhanced reproducibility.

\end{document}
